% \iffalse meta-comment
%
% hit-list.dtx 是各个类共用的列表环境
%
% \fi
%
% \section{共享列表环境}
%
% \changes{v3.2a}{2026/09/28}{列表照 \pkg{docxlike} 的多级列表排：编号格式、位置、编号之后跟什么都照 Word 的定义，
%   \pkg{enumitem} 的参数由它算}
% \changes{v3.2a}{2026/08/10}{\pkg{enumitem} 的加载、列表间距与 \env{publist} 抽出来共用，
%   \env{publist} 改用 \cs{newlist} 实现，报告也能写}
%
% 列表相关的东西都收在这里：\pkg{enumitem} 的加载、默认列表环境的全局间距、
% 以及 \env{publist}。
%
% \begin{macro}{\hit@setup@list@packages,\hit@setup@list@spacing}
% 分成两个函数是因为时机不同：宏包要早加载（\pkg{tocloft} 之类会跟它抢 \cs{cft*}
% 长度，顺序动不得），间距设置要晚一步。各个类在原来那两段代码所在的位置分别调用。
%
% 间距的取值不写死在这里，三个变量由各个类自己赋值：不赋就是 \hologo{LaTeX} 的默认间距，
% 赋了就按赋的来。这样三个类走同一条代码路径，「这个类要不要调整列表间距」是一次
% 显式的赋值，而不是靠某个类少调了一个函数来体现。
%    \begin{macrocode}
%<*shared-list>
\bool_new:N \g__hit_list_nosep_bool
\clist_new:N \g__hit_list_itemize_clist
\clist_new:N \g__hit_list_enumerate_clist
\cs_new_protected:Npn \hit@setup@list@packages
  { \RequirePackage{enumitem} }
\cs_new_protected:Npn \hit@setup@list@spacing
  {
    \bool_if:NT \g__hit_list_nosep_bool { \setlist { nosep } }
    \clist_if_empty:NF \g__hit_list_itemize_clist
      { \exp_args:NV \setitemize \g__hit_list_itemize_clist }
    \clist_if_empty:NF \g__hit_list_enumerate_clist
      { \exp_args:NV \setenumerate \g__hit_list_enumerate_clist }
  }
%    \end{macrocode}
% \end{macro}
%
% \begin{macro}{\hit@list@docxlike:nnn}
% \cs{hit@list@docxlike:nnn}\marg{\pkg{enumitem} 列表}\marg{多级列表}\marg{级数}：照 \pkg{docxlike} 的多级列表
% （Word“定义新多级列表”那一套，写在 \file{hithesis.cfg} 的 \option{multilevel-lists} 里）给 \pkg{enumitem}
% 的列表设参数。\pkg{enumitem} 只当排版器：标号是哪几个字、摆在哪、后面空多少、折行缩进到哪，都从
% Word 的定义算：
% \begin{itemize}
% \item 标号照这一级的“输入编号的格式”，\verb|\%|\meta{k} 换成第 \meta{k} 级的号、用第 \meta{k} 级的编号
%   格式；项目符号（\texttt{bullet}）照印。编号格式登记成 \pkg{enumitem} 的计数器写法
%   （\cs{AddEnumerateCounter}，与 \cs{arabic*} 同一路），这一级的号写成带星号的那种，上级的写明计数器，
%   标号与引用什么时候取值都交给 \pkg{enumitem}（它在列表开头展开一遍标号，那时计数器还是零，
%   直接读计数器就取成零）。装盒用
%   \cs{docxlike\_format\_run\_hbox\_set:Nn}，“（1）”这种头尾的全角括号占整格，与 Word 一样。
% \item 编号位置是 \texttt{ind/left} 加首行缩进（悬挂时减），就是首行起点；文本缩进位置 \texttt{ind/left} 是
%   \cs{leftmargin}，折行回到这里。
% \item 编号之后：空格是 Arial 的空格（\cs{c\_docxlike\_numbering\_space\_tl} 个字号），落成 \texttt{labelsep}；
%   什么都不跟是零；制表符跳到文本缩进位置（悬挂缩进的那个停靠点），落成标号列宽 \texttt{labelwidth}，
%   标号比它宽时照标号自己的宽。
% \end{itemize}
% \pkg{enumitem} 把标号摆在“\cs{leftmargin} 加 \cs{itemindent} 减 \cs{labelsep} 减 \cs{labelwidth}”处，
% 所以 \cs{itemindent} 取“编号位置减文本缩进位置加列宽加间隔”，首行正好从编号位置起。
% 列表的计数器名：\env{enumerate} 是 \texttt{enumi} 一族，\cs{newlist} 出来的是列表名加罗马数字。
%    \begin{macrocode}
\tl_new:N \l__hit_list_value_tl
\tl_new:N \l__hit_list_label_tl
\tl_new:N \l__hit_list_suff_tl
\dim_new:N \l__hit_list_left_dim
\dim_new:N \l__hit_list_first_dim
\dim_new:N \l__hit_list_width_dim
\tl_new:N \l__hit_list_sep_tl
\box_new:N \l__hit_list_label_box
\cs_new:Npn \__hit_list_counter:nn #1#2
  { \str_if_eq:nnTF {#1} { enumerate } { enum } {#1} \int_to_roman:n {#2} }
\cs_new_protected:Npn \__hit_list_twips:nnN #1#2#3
  {
    \docxlike_numbering_get:nnnN {#1} {#2} {#3} \l__hit_list_value_tl
    \tl_if_blank:VT \l__hit_list_value_tl { \tl_set:Nn \l__hit_list_value_tl { 0 } }
  }
\cs_new_protected:Npn \hit@list@label:n #1
  {
    \docxlike_format_run_hbox_set:Nn \l__hit_list_label_box {#1}
    \box_use_drop:N \l__hit_list_label_box
  }
\cs_new_protected:Npn \__hit_list_format:n #1
  {
    \cs_if_exist:cF { hitlist@ #1 }
      {
        \cs_gset:cpn { @hitlist@ #1 } ##1
          { \docxlike_number_format:nn {#1} { \int_use:N ##1 } }
        \cs_gset:cpn { hitlist@ #1 } ##1
          { \use:c { @hitlist@ #1 } { \use:c { c@ ##1 } } }
        \exp_args:Ncc \AddEnumerateCounter { hitlist@ #1 } { @hitlist@ #1 } { 0 }
      }
  }
\cs_new_protected:Npn \hit@list@docxlike:nnn #1#2#3
  {
    \int_step_inline:nn {#3}
      {
        \docxlike_numbering_get:nnnN {#2} {##1} { lvlText } \l__hit_list_label_tl
        \int_step_inline:nn { 9 }
          {
            \tl_if_in:NnT \l__hit_list_label_tl { \% ####1 }
              {
                \docxlike_numbering_get:nnnN {#2} {####1} { numFmt } \l__hit_list_value_tl
                \exp_args:NV \__hit_list_format:n \l__hit_list_value_tl
                \int_compare:nNnTF {####1} = {##1}
                  {
                    \tl_replace_all:Nne \l__hit_list_label_tl { \% ####1 }
                      { \exp_not:c { hitlist@ \l__hit_list_value_tl } * }
                  }
                  {
                    \tl_replace_all:Nne \l__hit_list_label_tl { \% ####1 }
                      {
                        \exp_not:c { hitlist@ \l__hit_list_value_tl }
                          { \__hit_list_counter:nn {#1} {####1} }
                      }
                  }
              }
          }
        \__hit_list_twips:nnN {#2} {##1} { ind/left }
        \dim_set:Nn \l__hit_list_left_dim { \l__hit_list_value_tl bp / 20 }
        \__hit_list_twips:nnN {#2} {##1} { ind/firstLine }
        \dim_set:Nn \l__hit_list_first_dim { \l__hit_list_value_tl bp / 20 }
        \__hit_list_twips:nnN {#2} {##1} { ind/hanging }
        \dim_sub:Nn \l__hit_list_first_dim { \l__hit_list_value_tl bp / 20 }
        \docxlike_numbering_get:nnnN {#2} {##1} { suff } \l__hit_list_suff_tl
        \dim_zero:N \l__hit_list_width_dim
        \str_case:VnF \l__hit_list_suff_tl
          {
            { space } { \tl_set:Ne \l__hit_list_sep_tl { \c_docxlike_numbering_space_tl em } }
            { nothing } { \tl_set:Nn \l__hit_list_sep_tl { 0pt } }
          }
          {
            \tl_set:Nn \l__hit_list_sep_tl { 0pt }
            \dim_set:Nn \l__hit_list_width_dim
              { \dim_max:nn { \c_zero_dim } { - \l__hit_list_first_dim } }
          }
        \use:e
          {
            \exp_not:N \setlist [ #1 , ##1 ]
              {
                leftmargin = \dim_use:N \l__hit_list_left_dim ,
                itemindent = \exp_not:N \dimexpr \dim_use:N \l__hit_list_first_dim
                  + \dim_use:N \l__hit_list_width_dim + \l__hit_list_sep_tl \exp_not:N \relax ,
                labelwidth = \dim_use:N \l__hit_list_width_dim ,
                labelsep = \l__hit_list_sep_tl ,
                align = left ,
                label = \exp_not:N \hit@list@label:n { \exp_not:V \l__hit_list_label_tl } ,
                ref = \exp_not:V \l__hit_list_label_tl ,
              }
          }
      }
  }
\cs_generate_variant:Nn \tl_replace_all:Nnn { Nne }
\cs_generate_variant:Nn \tl_if_blank:nT { V }
\cs_generate_variant:Nn \str_case:nnF { V }
%    \end{macrocode}
% \end{macro}
%
% \begin{macro}{\hit@list@word:nn}
% 照 Word 排的那几档（学位论文、照终稿排正文的报告）列表都是 Word 的多级列表：编号在首行缩进
% 处（正文的首行缩进，两个字），续行回到版心左缘，与范例里手敲“（1）……”的段落一个样。
% \cs{hit@list@word:nn}\marg{一级编号}\marg{编号之后} 定两个多级列表 \texttt{hit-enumerate} 与
% \texttt{hit-itemize}，再照它们设 \env{enumerate} 与 \env{itemize}。一级编号用“输入编号的格式”
% 的写法（\verb|（\%1）|、\verb|\%1.|），编号之后取 \texttt{space} 或 \texttt{nothing}；二到四级照
% \pkg{enumitem} 原来的标号（(a)、i.、A.；项目符号 \textbullet、\textendash、\textasteriskcentered、
% \textperiodcentered），后面跟空格，一级套一级、各自从上一级的正文起点量。
%
% 首行缩进要等正文的样式设好才有值，所以在 \texttt{begindocument/end} 上定，从 \texttt{Normal}
% 现算。\pkg{enumitem} 的 \cs{setlist} 在正文里也能写。
%    \begin{macrocode}
\cs_new_protected:Npn \hit@list@word:nn #1#2
  {
    \AddToHook { begindocument / end }
      {
        \group_begin:
          \docxlike_format_compute:n { Normal }
          \dim_gset_eq:NN \g__hit_list_indent_dim \l_docxlike_format_indent_dim
        \group_end:
        \use:e
          {
            \exp_not:N \docxlikesetup
              {
                multilevel-lists =
                  {
                    hit-enumerate =
                      {
                        level-1 = { \__hit_list_level:nnn { \exp_not:n {#1} } { decimal } {#2} } ,
                        level-2 = { \__hit_list_level:nnn { ( \exp_not:N \% 2 ) } { lower-letter } { space } } ,
                        level-3 = { \__hit_list_level:nnn { \exp_not:N \% 3 . } { lower-roman } { space } } ,
                        level-4 = { \__hit_list_level:nnn { \exp_not:N \% 4 . } { upper-letter } { space } } ,
                      } ,
                    hit-itemize =
                      {
                        level-1 = { \__hit_list_level:nnn { \exp_not:N \textbullet } { bullet } { space } } ,
                        level-2 = { \__hit_list_level:nnn { \exp_not:N \textendash } { bullet } { space } } ,
                        level-3 = { \__hit_list_level:nnn { \exp_not:N \textasteriskcentered } { bullet } { space } } ,
                        level-4 = { \__hit_list_level:nnn { \exp_not:N \textperiodcentered } { bullet } { space } } ,
                      } ,
                  }
              }
          }
        \hit@list@docxlike:nnn { enumerate } { hit-enumerate } { 4 }
        \hit@list@docxlike:nnn { itemize } { hit-itemize } { 4 }
        \setlist { nosep }
      }
  }
\dim_new:N \g__hit_list_indent_dim
\cs_new:Npn \__hit_list_level:nnn #1#2#3
  {
    enter-formatting-for-number = {#1} ,
    number-style-for-this-level = #2 ,
    aligned-at = \dim_use:N \g__hit_list_indent_dim ,
    text-indent-at = 0pt ,
    follow-number-with = #3
  }
%    \end{macrocode}
% \end{macro}
%
% \DescribeEnv{publist} 发表文章、专利、科研项目的编号列表，用法是
% \cs{begin}\marg{publist} 里逐条 \cs{item}。
% 参数照搬原来那个 \env{list} 的设置。计数器由 \pkg{enumitem} 内部管，模板这边
% 就不必再造一个 LaTeX 计数器了。\cs{labelwidth} 原来没显式设过，靠的是
% \cs{leftmargin}\texttt{-}\cs{itemindent}\texttt{-}\cs{labelsep}，这里写明。
%
% \begin{macro}{\hit@setup@publist}
% 标号的排法照搬内核 \env{list} 的 \cs{@mklab}：装进一个 \cs{labelwidth} 宽的盒子、
% 用 \cs{hfil} 顶到右边。\pkg{enumitem} 自带的 \texttt{align=right} 用的是 \cs{hss}，
% 标号超宽时会往左溢出，与原来相反。
%    \begin{macrocode}
\cs_new_protected:Npn \hit@setup@publist
  {
\SetLabelAlign{hit@pubright}{\hb@xt@\labelwidth{\hfil##1}}
\newlist{publist}{enumerate}{1}
\setlist[publist]{%
  label={[\arabic*]},
  leftmargin=1.7em,
  itemindent=0em,
  labelsep=0.5em,
  labelwidth=24pt,
  align=hit@pubright,
  rightmargin=0em,
  topsep=0ex,
  parsep=0ex,
  itemsep=0ex,
  listparindent=0pt,
  partopsep=3pt plus 1pt minus 1pt,
}
  }
%    \end{macrocode}
% \end{macro}
%
%    \begin{macrocode}
%</shared-list>
%    \end{macrocode}
%
% \Finale
